% Part: computability
% Chapter: computability-theory
% Section: universal-part-function

\documentclass[../../../include/open-logic-section]{subfiles}

\begin{document}

\olfileid{cmp}{thy}{uni}
\olsection{تابع محاسبه‌پذیر جزئیِ جهانی}

\begin{thm}
\ollabel{thm:univ-comp}
تابع محاسبه‌پذیر جزئیِ جهانی~$\fn{Un}(e,x)$ وجود دارد. به بیان دیگر،
تابعی چون~$\fn{Un}(e,x)$ با ویژگی‌های زیر وجود دارد:
\begin{enumerate}
\item $\fn{Un}(e,x)$ محاسبه‌پذیر جزئی است.
\item اگر $f(x)$ هر تابع محاسبه‌پذیر جزئی باشد، عدد طبیعی~$e$ای وجود
دارد چنان‌که برای هر~$x$،
$f(x) \simeq \fn{Un}(e,x)$.
\end{enumerate}
\end{thm}

\begin{proof}
قرار دهید $\fn{Un}(e,x) \simeq U(\umin{s}{T(e,x,s)})$، که در آن
$U$ و~$T$ همان‌هایی هستند که در قضیهٔ صورت نرمال کلینی آمده‌اند
(\olref[nfm]{thm:normal-form}).
\end{proof}

\begin{explain}
این تنها راهی دقیق برای گفتن این است که برشماری مؤثری از توابع
محاسبه‌پذیر جزئی داریم. اندیشه این است که اگر~$f_e$ را برای تابعی
بنویسیم که با $f_e(x)=\fn{Un}(e,x)$ تعریف می‌شود، آنگاه دنبالهٔ
$f_0$، $f_1$، $f_2$، \dots همهٔ توابع محاسبه‌پذیر جزئی را دربر
می‌گیرد و~$f_e(x)$ را می‌توان به‌طور «یکنواخت» برحسب~$e$ و~$x$
محاسبه کرد. برای سادگی از تابعی دوموضعی استفاده می‌کنیم که برای توابع
یک‌موضعی جهانی است؛ اما با کدگذاری دنباله‌های اعداد می‌توان این مطلب
را به‌آسانی به آرگومان‌های بیشتر تعمیم داد. برای نمونه، توجه کنید اگر
$f(x,y,z)$ تابعی بازگشتی جزئی و سه‌موضعی باشد، آنگاه
$g(x) \simeq f((x)_0,(x)_1,(x)_2)$ تابعی بازگشتی جزئی و یک‌موضعی است.
\end{explain}

\end{document}
